\input _template

\setlanguage{EN}

\Title{Algebra}

\Subtitle{Factorization}

\MathcompsLink{rozklady-na-sucin}

\Author{Patrik Bak}

\sec Introduction

Why is it good to know how to factor complex expressions? Imagine you have to prove that the number $n^3-n$ is divisible by six for every integer $n$. Without modification, it is unclear. However, it is enough to factor the expression:
$$
n^3-n = n(n^2-1) = (n-1)n(n+1).
$$
Suddenly we see a product of three consecutive numbers, of which at least one is always even and exactly one is divisible by three. Divisibility by six is thus obvious.

The ability to transform a confusing sum into a clear product is one of the key techniques in solving tricky mathematical problems\fnote{This idea of factoring expressions is thousands of years old -- ancient Babylonians and Greeks solved problems, which we record today as quadratic equations, using geometric methods that were essentially a visual factorization.}. In this lesson, we will systematically derive known factorizations, explain the tricks and intuition behind them, and practice everything on examples.

\sec Theory 

In this section, we will gradually introduce four key factorization techniques. We will start from basic formulas, move to the general method of grouping, and finally explain completing the square. As a curiosity, we will also show that one can complete not only to a square but also to a cube, and surprisingly it can be useful.

\secc Difference and sum of powers formulas

\Exercise{}{
    Verify by expanding that the following holds
    \begitems \style i
    \i $a^2-b^2=(a-b)(a+b)$
    \i $a^3-b^3=(a-b)(a^2+ab+b^2)$
    \i $a^4-b^4=(a-b)(a^3+a^2b+ab^2+b^3)$
    \enditems
    What will the general formula for $a^n-b^n$ look like for a natural number $n$?
}{
    The answer to the question is Proposition 1.
}

\Exercise{}{
    Verify by expanding that the following holds
    \begitems \style i
    \i $a^3+b^3=(a+b)(a^2-ab+b^2)$
    \i $a^5+b^5=(a+b)(a^4-a^3b+a^2b^2-ab^3+b^4)$
    \enditems
    What will the general formula for $a^n+b^n$ look like for an odd natural number $n$?
}{
    The answer to the question is Proposition 2.
}

After a while of investigation, we will surely come up with general formulas:

\Theorem{}{
    For all real numbers $a,b$ and natural $n$, it holds
    $$
    a^n - b^n = (a-b)(a^{n-1}+a^{n-2}b+\cdots+ab^{n-2}+b^{n-1})
    $$
}{
    We subtract the expressions:
    $$
    \align
    a(a^{n-1}+a^{n-2}b+\cdots+ab^{n-2}+b^{n-1}) &= a^n  + a^{n-1}b + \cdots + ab^{n-1}  \\
    b(a^{n-1}+a^{n-2}b+\cdots+ab^{n-2}+b^{n-1}) &= \hskip1cm a^{n-1}b + \cdots + ab^{n-1} + b^n
    \endalign
    $$
    On the left, we can factor out the large parenthesis and get exactly the right side of the proven formula. On the right, all terms cancel out. The proposition is proved.
}

The situation with $a^n+b^n$ is more curious in that such a factorization works only for odd~$n$. From the proof, it is visible why:

\Theorem{}{
    For all real numbers $a,b$ and odd natural $n$, it holds
    $$
    a^n + b^n = (a+b)(a^{n-1}-a^{n-2}b+\cdots-ab^{n-2}+b^{n-1})
    $$
}{
    The proposition can be proved similarly to the previous one -- by observing which terms cancel out when expanding the right side. However, a trickier proof is worth mentioning, in which we substitute the value $-b$ instead of $b$ into the formula for $a^n-b^n$. Thanks to the oddness of $n$, then $a^n-(-b)^n=a^n+b^n$. Every second term of the parenthesis $a^{n-1}+a^{n-2}b+\cdots+ab^{n-2}+b^{n-1}$ changes sign when replacing $b$ with $-b$.
}

How is it for even $n$? For the expression $a^2+b^2$, it can be proven that it really cannot be factored (in the domain of real numbers). However, all other $a^n+b^n$ can be factored, which we will gradually work towards. 

\Exercise{}{
    Factor into a product of two expressions without roots ($n,k$ are natural numbers).
    \begitems \style i
    \i $a^3-27$
    \i $a^3+8b^3$
    \i $a^4-4^n$
    \i $a^6+b^6$
    \i $a^{4k+2}+b^{4k+2}$
    \enditems
}{
    The individual factorizations are
    \begitems \style i
    \i $a^3-27 = a^3 - 3^3 = (a-3)(a^2+3a+9)$
    \i $a^3+8b^3 = a^3 + (2b)^3 = (a+2b)(a^2-2ab+4b^2)$
    \i $a^4-4^n = (a^2)^2 - (2^n)^2 = (a^2-2^n)(a^2+2^n)$
    \i $a^6+b^6 = (a^2)^3 + (b^2)^3 = (a^2+b^2)(a^4 - a^2b^2 + b^4)$
    \i $a^{4k+2}+b^{4k+2}=(a^{2})^{2k+1}+(b^{2})^{2k+1}=(a^{2}+b^{2})(a^{4k}-a^{4k-2}b^{2}+\cdots-a^{2}b^{4k-2}+b^{4k})$
    \enditems
}

In more complex problems, factorization is typically just one of several steps. We can try:

\Problem{0}{}{
    Prove that for every natural number~$n$, the number $n^5-n$ is divisible by 30.
}{
    Factor the investigated expression into as many parentheses as possible.
}{
    The sought factorization is $n^5-n=n(n^4-1)=n(n^2-1)(n^2+1)=n(n-1)(n+1)(n^2+1)$. To investigate divisibility by 30, it is enough to investigate divisibility by $2$, $3$, $5$ separately.
}{
    It holds that $n^5-n=n(n^4-1)=n(n^2-1)(n^2+1)=n(n-1)(n+1)(n^2+1)$. It suffices to prove that this expression is divisible by $2$, $3$, $5$. 
    \begitems
    \style i
    \i divisibility by $2$ follows from the fact that of consecutive numbers $n$, $n-1$, at least one is even;
    \i similarly, divisibility by $3$ follows from numbers $n-1$, $n$, $n+1$;
    \i divisibility by $5$ is more complex. Certainly, when $n$ gives a remainder of 0, 1, or 4 when divided by 5, we are done, due to factors $n$, $n-1$, $n+1$. Then if $n=5k+2$ or $n=5k+3$, then $n^2+1 = 25k^2+20k+5$ or $n^2+1=25k^2+30k+10$, so again we have a number divisible by 5.
    \enditems
}

\Problem{0}{}{
    Find three distinct pairs of natural numbers $(a,b)$ such that the number $a^{37} + b^{37}$ is divisible by~$7$ and~$0<a<b<7$.
}{
    The factorization of $a^{37}+b^{37}$ suggests a much simpler divisibility, which is sufficient.
}{
    Since $a^{37}+b^{37}=(a+b)(\cdots)$, it is enough that $7 \mid a+b$.
}{
    It holds that $a^{37}+b^{37}=(a+b)(a^{36}+\cdots+b^{36})$, so $a+b \mid a^{37}+b^{37}$, thus it suffices that $7 \mid a+b$. We can easily achieve this with pairs $(a,b)$ equal to $(1,6), (2,5), (3,4)$.
}

\Problem{1}{}{
    Find all natural numbers $n$ greater than 1 such that $n^6-1$ is a product of three not necessarily distinct primes.
}{
    The expression $n^6-1$ can be factored into a product of many parentheses.
}{
    One possible factorization is $n^6-1=(n^2)^3-1=(n^2-1)(n^4-n^2+1)$. The second parenthesis can indeed be factored, but it is not evident at all how. A better factorization is $n^6-1=(n^3)^2-1 = (n^3-1)(n^3+1)$. From this we factor further. Subsequently, we have at least 4 factors, which sounds suspicious if the number is to be a product of three primes.
}{
    We have the factorization
    $$
    n^6-1=(n^3)^2-1 = (n^3-1)(n^3+1)=(n-1)(n^2+n+1)(n+1)(n^2-n+1).
    $$
    For this to be a product of three primes, one parenthesis must be equal to 1. For $n>1$, this is only possible for the first one, where $n=2$. In that case, we indeed have $2^6-1=63=3\cdot3\cdot7$.
}

\secc Grouping

The common school procedure of successive factorization works even in harder problems. The idea is: notice that something can be factored out; factor it out; and see what happens next. For this procedure, there exists a very important trick usable even in harder problems: \textit{watch when the expression is zero}.

\Example{}{
    Factor the expression $a^4-a-b^4+b$. 
}{
    Without seeing the resulting factorization, we can say that $a-b$ will be in it, because the investigated expression is equal to 0 for $a=b$. Thanks to this, we purposefully rearrange the terms as $(a^4-b^4)+(a-b)$. Now we can factor out $a-b$ from both expressions and we have 
    $$
    (a-b)(a^3+a^2b+ab^2+b^3)-(a-b)=(a-b)(a^3+a^2b+ab^2+b^3-1).
    $$
}

This view explains why in the formulas $a^n-b^n$ and $a^n+b^n$ from the previous section we have $a-b$ and $a+b$; and also why the second one requires odd $n$ -- for $a=b$ or $a=-b$, $a^n-b^n$ or $a^n + b^n$ are zero.

\Exercise{}{Factor: 
\begitems \style i
\i [2 factors] $2ab+a+2b+1$
\i [3 factors] $abc+ab+bc+ca+a+b+c+1$
\i [3 factors] $a^2(b-c) + b^2(c-a) + c^2(a-b)$
\i [4 factors] (harder\fnote{Coming up with this factorization was the first step to solving problem 3 of IMO 2006}) $ab(a^{2}-b^{2})+bc(b^{2}-c^{2})+ca(c^{2}-a^{2})$
\enditems}{
    The individual factorizations are:
    \begitems \style i
    \i An obvious zero point is $a=-1$, in the result we expect $a+1$. We easily find $2ab+a+2b+1=(a+1)(2b+1)$.
    \i Here again it holds that whenever any of the numbers $a,b,c$ is equal to $-1$, the expression will be zero. Thus we expect $(a+1)(b+1)(c+1)$ in the result. Actually, that is all: $abc+ab+bc+ca+a+b+c+1=(a+1)(b+1)(c+1)$. We can arrive at this gradually like this: 
    $$
    \gather
    abc+ab+bc+ca+a+b+c+1=
    (abc+ab)+(bc+b)+(ac+a)+(c+1) = \\ =
    ab(c+1)+b(c+1)+a(c+1)+(c+1)=
    (c+1)(ab+b+a+1)=  \\ =
    (c+1)(b(a+1)+(a+1))=
    (c+1)(b+1)(a+1).
    \endgather
    $$
    \i When two numbers are equal, e.g., $a=b$, the expression is zero. We manipulate the expression to have a parenthesis $a-b$ in the result. We expect that $b-c$, $c-a$ (or the opposite) will also appear there. 
    $$ 
    \gather
    a^{2}(b-c)+b^{2}(c-a)+c^{2}(a-b) = (a^2b - ab^2) - c(a^2-b^2) + c^2(a-b) = \\ =
    ab(a-b) - c(a-b)(a+b) - c^2(a-b) = (a-b)(ab-c(a+b)-c^2) = \\ = (a-b)(a(b-c)-c(b-c)) = (a-b)(b-c)(a-c).
    \endgather
    $$
    \i We proceed similarly to the previous exercise, since again we have equality for $a=b$. The modifications are more complex here:
    $$
    \gather
    ab(a^{2}-b^{2})+bc(b^{2}-c^{2})+ca(c^{2}-a^{2})=
    ab(a^2-b^2) + b^3c - bc^3 + ac^3 - a^3c = \\ =
    ab(a^2-b^2) - c(a^3-b^3) + c^3(a-b) = 
    (a-b)(ab(a+b) - c(a^2+ab+b^2) + c^3).
    \endgather
    $$
    Now let's focus on the parenthesis:
    $$
    \gather
    ab(a+b) - c(a^2+ab+b^2) + c^3 = 
    a^2b + ab^2 - ca^2 - abc - cb^2 + c^3 = \\ =
    ab(a-c) + b^2(a-c) - c(a^2-c^2) = 
    (a-c)(ab+b^2-c(a+c)).
    \endgather
    $$
    Finally, the last parenthesis:
    $$
    \gather
    ab+b^2-c(a+c)=(ab-ac)+(b^2-c^2)=\\=a(b-c) + (b-c)(b+c)=  (b-c)(a+b+c).
    \endgather
    $$
    Together we have 
    $$
    a^{2}(b-c)+b^{2}(c-a)+c^{2}(a-b) = (a-b)(b-c)(a-c)(a+b+c).
    $$
    \enditems
}

This view is not universal, because the factored factors do not have to be linear at all or can be more complex and it is harder to see the root. Then we have no choice but to play with rearranging terms and successive factoring. 

\Exercise{}{
    Factor into a product of two expressions: 
    \begitems \style i
    \i $a^3b + a^2 + ab^3 + b^2$
    \i $a^3b + a^2 + ab^3 + ab + b^2 + 1$
    \i $a^3+b^3+c^3+ab(a+b)+bc(b+c)+ca(c+a)$
    \enditems
}
{
    The individual factorizations are:
    \begitems \style i
    \i $a^3b + a^2 + ab^3 + b^2 = a^2(ab+1) + b^2(ab+1)=(a^2+b^2)(ab+1)$.
    \i
    $$
    \abovedisplayskip-19pt
    \gather
    a^3b + a^2 + ab^3 + ab + b^2 + 1 = (a^3b+a^2) + (ab^3+b^2) + (ab+1) =\\ 
    = a^2 (ab+1) + b^2 (ab+1) + (ab+1) = (a^2+b^2+1)(ab+1).
    \endgather
    $$
    \i
    $$  \abovedisplayskip-19pt
    \gather
    a^3+b^3+c^3+ab(a+b)+bc(b+c)+ca(c+a)= \\=
    (a^3+ab^2+ac^2) + (b^3+bc^2+ba^2) + (c^3+ca^2+cb^2) = \\ =
    a(a^2+b^2+c^2) + b(a^2+b^2+c^2) + c(a^2+b^2+c^2) = \\ =
    (a+b+c)(a^2+b^2+c^2).
    \endgather
    $$
    \enditems
}

Hidden factorization and other algebraic modifications can be successfully used in this example:

\Problem{1}{CPSJ\fnote{Czech-Polish-Slovak Junior Match} 2018}{For natural numbers $a,b,c$ it holds 
$$
(a + b + c)^2 \mid ab(a + b) + bc(b + c) + ca(c + a) + 3abc.
$$
Prove that
$$
(a + b + c) \mid (a - b)^2 + (b - c)^2 + (c - a)^2.
$$}{
    The expression $ab(a+b) + bc(b+c) + ca(c+a) + 3abc$ can be subtly factored.
}{
    It holds that $ab(a+b) + bc(b+c) + ca(c+a) + 3abc = (a+b+c)(ab+bc+ca)$. Thus our divisibility simplifies to
    $$
    a+b+c \mid ab+bc+ca.
    $$ 
    The expression $(a-b)^2 + (b-c)^2 + (c-a)^2$ is equal to $2(a^2+b^2+c^2-ab-bc-ca)$, we see a connection.
}{
    We factor the right side of the first divisibility:
    $$
    \align
    ab(a + b) + bc(b + c) + ca(c + a) + 3abc &= \\
    (ab(a + b)+abc) + (bc(b + c)+abc) + (ca(c + a)+abc) &= \\
    ab(a+b+c) + bc(a+b+c) + ca(a+b+c) &= \\
    (a+b+c)(ab+bc+ca).
    \endalign
    $$
    The assumed divisibility translates to $a+b+c \mid ab+bc+ca$. It holds 
    $$
    (a-b)^2 + (b-c)^2 + (c-a)^2 = 2(a^2+b^2+c^2-ab-bc-ca).
    $$
    The divisibility $a+b+c \mid ab+bc+ca$ simplifies the situation to the fact that it suffices to prove $a+b+c \mid 2(a^2+b^2+c^2)$. We prove this by using
    $$
    a+b+c \mid (a+b+c)^2 = a^2+b^2+c^2 + 2(ab+bc+ca).
    $$
    Together with $a+b+c \mid ab+bc+ca$, we then have $a+b+c \mid a^2+b^2+c^2$, so altogether we are done.
}

\secc Completing the square

Completing the square is one of the oldest and most elegant techniques in algebra\fnote{It was made famous by the Persian mathematician of the 9th century, Muhammad ibn Musa al-Khwarizmi (the word 'algorithm' comes from his name). He did not solve equations with our symbols, but geometrically -- literally drawing squares and rectangles and trying to 'complete' a larger square from them.}. It allows us to \textit{factor} a quadratic polynomial into a product of simpler linear ones, e.g.
$$
x^2+2x-3=(x+1)^2-4=(x+1-2)(x+1+2)=(x-1)(x+3).
$$
In general for real numbers $a\neq 0$, $b$, $c$ we have
$$
\gather
ax^2+bx+c
=a\(x^2+\frac ba x + \frac ca\)
=a\(\(x + \frac{b}{2a}\)^2 + \frac ca - \frac{b^2}{4a^2}\) =\\
=a\(\(x + \frac{b}{2a}\)^2 - \frac{b^2 - 4ac}{4a^2}\).
\endgather
$$
From this, we can easily derive the formula for the solutions of a quadratic equation.

In general, we actually used $a^2+2ab = (a+b)^2 - b^2$. However, we can also complete to a square by 'manufacturing' the middle coefficient, i.e., using the formula $a^2+b^2=(a+b)^2-2ab$. We will demonstrate this on the following example:

\Example{}{
    Factor $m^4+m^2n^2+n^4$.
}{
    The first thing that might occur to us is to complete $m^4+m^2n^2$ to a square. By doing so we get 
    $$
    m^4+m^2n^2+n^4= \(m^2 +\frac{n^2}2\)^2 + \frac{3n^4}4.
    $$
    This did not help us much. Let's try something else, let's complete $m^4+n^4$ to a square. Then we have 
    $$
    \gather
    m^4+m^2n^2+n^4 = (m^4+n^4)+m^2n^2 = (m^2+n^2)^2 - 2m^2n^2 + m^2n^2 = \\ = (m^2+n^2)^2 - m^2n^2 = (m^2+n^2-mn)(m^2+n^2+mn).
    \endgather
    $$
}

Practice this by deriving a known identity:

\Exercise{Sophie-Germain identity\fnote{Named after Sophie Germain, one of the most significant female mathematicians at the turn of the 18th and 19th centuries. At a time when women did not have access to universities, she studied mathematics herself and corresponded with the greatest mathematicians of her time, like Gauss, initially under a male pseudonym.}
}{
    Factor $a^4+4b^4$.
}{
    By completing to a square we have 
    $$
    a^4+4b^4 = (a^2)^2 + (2b^2)^2 = (a^2+2b^2)^2 - (2ab)^2 = (a^2+2b^2-2ab)(a^2+2b^2+2ab).
    $$
}

Try an example from the national round of MO:

\Problem{0}{CKMO 2012}{
    Determine all natural numbers $n$ for which $n^4 - 3n^2 + 9$ is a prime number.
}{
    The investigated expression can surprisingly be factored. Completing the square is a fine technique. Just be careful what we are completing.
}{
    The key completion is 
    $$
    n^4 - 3n^2 + 9 = (n^4 + 9) - 3n^2 = (n^2+3)^2 - 6n^2 - 3n^2=(n^2+3)^2-9n^2.
    $$
    Do we see a known formula?
}{
    It holds 
    $$
    \gather
    n^4 - 3n^2 + 9 = (n^4 + 9) - 3n^2 = (n^2+3)^2 - 6n^2 - 3n^2=\\=(n^2+3)^2-9n^2 = (n^2+3-3n)(n^2+3+3n).
    \endgather
    $$
    Clearly $n^2+3+3n > n^2+3-3n$. The second number is positive for $n>0$, because 
    $$
    n^2-3n+3=\(n-\frac32\)^2 + \frac 34.
    $$
    For the product to be a prime number, since $(n^2+3-3n)(n^2+3+3n)$, then $n^2+3-3n=1$. By solving the quadratic equation we get $n=1$ and $n=2$. 
}

Learned techniques can be nicely put together in this observation:

\Problem{1}{}{
    Justify that the expression $a^n+b^n$ can be factored into a product of non-constant expressions with real coefficients for every $n>2$.
    
    (Let us note that it can be proven that $a^2+b^2$ \textit{cannot} be factored into a product of two non-constant parentheses with real coefficients.)
    }{
    For odd $n$ we have a formula. Let's realize that it is even enough for $n$ to have an odd divisor. Thanks to this observation, the even case is also simplifying. Can we solve the first even case $n=4$? Does it help with the rest?
}{
    If $n$ has an odd divisor $k$, where $n=mk$, then we can factor out $a^k+b^k$ from $a^n+b^n$. It remains to solve numbers that do not have an odd divisor -- that is, powers of 2. According to the problem statement $n>2$, so the smallest interesting power is 4, for such powers we have completing the square. But what about higher powers? Well, those are fortunately divisible by~4.
}{
    If $n$ has an odd divisor $k$, where $n=mk$, then
    $$
    \gather
    a^n+b^n
        = a^{mk}+b^{mk}
        = (a^m)^k+(b^m)^k =\\
        = (a^m+b^m)\(a^{m(k-1)}-a^{m(k-2)}b^{m}+a^{m(k-3)}b^{2m}-\cdots
        - a^{m}b^{m(k-2)}+b^{m(k-1)}\).
    \endgather
    $$
    Otherwise, $n$ is a power of~2 and thanks to $n>2$ it is divisible by 4, so $n=4t$. Then 
    $$
    \gather
    a^{4t}+b^{4t}
        =(a^{2t})^{2}+(b^{2t})^{2}
        =(a^{2t}+b^{2t})^{2}-2a^{2t}b^{2t}=\\
        =\(a^{2t}+b^{2t}-\sqrt2\,a^{t}b^{t}\)\(a^{2t}+b^{2t}+\sqrt2\,a^{t}b^{t}\)=\\
        =\(a^{2t}-\sqrt2\,a^{t}b^{t}+b^{2t}\)\(a^{2t}+\sqrt2\,a^{t}b^{t}+b^{2t}\).
    \endgather
    $$ 
}

A few more harder problems for practice:

\Problem{1}{}{
    Factor $2a^2b^2+2b^2c^2+2c^2a^2-a^4-b^4-c^4$ into a product of four non-constant parentheses.
}{
    Notice that $2a^2b^2-a^4-b^4$ appears in the expression, which is equal to $-(a^2-b^2)^2$. The key trick is to combine it with the rest of the expression using *another* completion to a square, where one of these 'squares' will be $-(a^2-b^2)^2$.
}{
    The second sought 'square' will be $-c^4$, as it holds: $-(a^2-b^2)^2 - c^4 = -(a^2-b^2+c^2)^2 + 2(a^2-b^2)c^2$. That fits quite well with the rest of the expression.
}{
    First we notice $2a^2b^2-a^4-b^4$ in the expression and create a square:
    $$
    2a^2b^2+2b^2c^2+2c^2a^2-a^4-b^4-c^4 = -(a^2-b^2)^2 - c^4 + 2b^2c^2 + 2c^2a^2.
    $$
    Now we apply completing the square to the first two terms, specifically we complete $-(x^2+y^2)$ for $x=a^2-b^2$ and $y=c^2$:
    $$
    \gather
    -(a^2-b^2)^2 - c^4 + 2b^2c^2 + 2c^2a^2= \\
    = -((a^2-b^2)+c^2)^2 + 2(a^2-b^2)c^2 + 2b^2c^2 + 2c^2a^2= \\
    = -(a^2-b^2+c^2)^2 + (2a^2c^2 - 2b^2c^2) + 2b^2c^2 + 2c^2a^2= \\
    = -(a^2-b^2+c^2)^2 + 4a^2c^2= \\
    = (2ac)^2 - (a^2-b^2+c^2)^2.
    \endgather
    $$
    We can easily factor this expression by repeatedly finding squares and using the formula for the difference of squares:
    $$
    \gather
    (2ac - (a^2-b^2+c^2))(2ac + (a^2-b^2+c^2)) = \\
    = (b^2 - (a^2-2ac+c^2))((a^2+2ac+c^2)-b^2)= \\
    = (b^2 - (a-c)^2)((a+c)^2 - b^2)= \\
    = (b-(a-c))(b+(a-c)) \cdot ((a+c)-b)((a+c)+b)= \\
    = (a+b+c)(a+b-c)(a-b+c)(-a+b+c).
    \endgather
    $$
    Let us note that we actually practically reconstructed Heron's\fnote{Heron of Alexandria was an ancient Greek mathematician and engineer who lived in the 1st century AD. His formula for the area of a triangle appears in his work \textit{Metrica}} formula. For the area~$S$ of a triangle with sides $a,b,c$, it holds 
    $$
    16S^2 = 2a^2b^2+2b^2c^2+2c^2a^2-a^4-b^4-c^4.
    $$
    This formula is usually presented in a computationally more acceptable~form 
    $$
    S = \sqrt{s(s-a)(s-b)(s-c)}, \qquad\text{where}\qquad s=\frac{a+b+c}{2}.
    $$
    Verify yourselves that this form is equivalent to our proven identity
    $$
    16S^2=2a^2b^2+2b^2c^2+2c^2a^2-a^4-b^4-c^4=(a+b+c)(a+b-c)(a-b+c)(-a+b+c).
    $$
}

\Problem{1}{}{
    For which natural numbers $n$ is the number $n^4 + 4^n$ a prime number?
}{
    The investigated expression resembles the Sophie-Germain identity. Try to fit it in there. Maybe it will be necessary to analyze some cases.
}{
    To be able to use Sophie-Germain, we need $4^n$ to be in the form $4b^4$. This is certainly true for odd $n$, because then
    $$
    4^n = 4 \cdot 4^{n-1} = 4 \cdot 2^{2(n-1)} = 4 \cdot \left(2^{\frac{n-1}{2}}\right)^4.
    $$  
    On the other hand, the case of even $n$ is again evidently non-prime.
}{
    We will analyze two cases according to the parity of $n$.

    \begitems \style i
    \i If $n$ is even, then the number $n^4+4^n$ is clearly divisible by 4, so it is not a prime number.
    \i If $n$ is odd, then notice that:
    $$
    4^n = 4 \cdot 4^{n-1} = 4 \cdot 2^{2(n-1)} = 4 \cdot \left(2^{(n-1)/2}\right)^4.
    $$  
    Our expression thus has the form $a^4+4b^4$ for $a=n$, $b=2^{\frac{n-1}2}$. Thus we can use the Sophie-Germain identity 
    $$
    a^4+4b^4=(a^2+2b^2-2ab)(a^2+2b^2+2ab).
    $$
    For this product to be a prime number, one of the factors must be equal to 1. The second factor is clearly greater than 1 for every $a,b \ge 1$. Let's analyze the first factor, which after completing to a square is equal to:
    $$
    a^2+2b^2-2ab=(a-b)^2+b^2.
    $$
    For this to be equal to 1, we must have $a=b=1$, so $n=1$. Then indeed $n^4+4^n=5$ is a prime number.
    \enditems

    The only solution is therefore $n=1$.
}

\secc Completing the cube

Finally, we will show a non-traditional technique: besides completing to a square, one can perform \textit{completing to a cube}. This means that if we have $a^3+b^3$, instead of directly applying the factorization formula, we can use 
$$
a^3 + b^3 = (a+b)^3 - 3ab(a+b).
$$
We will demonstrate this on an example:

\Example{}{
    Factor $a^3+b^3+c^3-3abc$.
}{
    Using completing to a cube we have
    $$
    \align
    a^3 + b^3 + c^3 - 3abc &= \\
    (a+b)^3 - 3ab(a+b) + c^3 - 3abc &= \\
    (a+b)^3 + c^3 - 3ab(a+b+c) &= \\
    (a+b+c)((a+b)^2 - (a+b)c + c^2) - 3ab(a+b+c) &= \\
    (a+b+c)((a+b)^2 - (a+b)c + c^2 - 3ab) &= \\
    (a+b+c)(a^2 + 2ab + b^2 - ac - bc + c^2 - 3ab) &= \\
    (a+b+c)(a^2 + b^2 + c^2 - ab - bc - ca).
    \endalign
    $$
}

This factorization is very interesting in itself, as it allows us to easily prove an inequality:

\Theorem{}{
    Let $a,b,c$ be real numbers for which $a+b+c \ge 0$. Then 
    $$
    a^3+b^3+c^3 \ge 3abc,
    $$
    where equality occurs if and only if $a+b+c=0$ or $a=b=c$.
}{
    We will use our already known factorization
    $$
    a^3+b^3+c^3 - 3abc = (a+b+c)(a^2+b^2+c^2-ab-bc-ca).
    $$
    Since $a+b+c \ge 0$, it suffices to prove $a^2+b^2+c^2 \ge ab+bc+ca$. This looks tempting from the perspective of completing to a square, for example, we can try to complete $a^2-ab$ to a square. However, I will reveal that this will not lead to the goal. The trick to proving this inequality is to multiply it by two and prove the equivalent inequality
    $$
    2a^2+2b^2+2c^2 \ge 2ab+2bc+2ca.
    $$
    It seems we haven't helped ourselves. However, another trick is to split $2a^2$ as $a^2+a^2$. After suitable rearranging of terms we then have 
    $$
    \align
    a^2+a^2+b^2+b^2+c^2+c^2-2ab-2bc-2ca &= \\
    (a^2-2ab+b^2)+(b^2-2bc+c^2)+(c^2-2ac+a^2) &=\\
    (a-b)^2 + (b-c)^2 + (c-a)^2.
    \endalign
    $$
    Since this last expression is certainly non-negative, the proof of the inequality is finished. Equality occurs if and only if the first parenthesis $a+b+c$ is zero or the second parenthesis $a^2+b^2+c^2-ab-bc-ca$ is zero. From the proof of its non-negativity, we see that it is exactly when $a-b=0$, $b-c=0$, $c-a=0$, i.e., when $a=b=c$. We are done.
}

\sec What to remember

\secc Techniques

\begitems \style N
\i We look for differences/sums of powers so we can use $a^n+b^n$ and $a^n-b^n$
\i We check when the expression is zero, which helps us factor correctly
\i We complete to a square/cube. We try multiple possibilities of how to do it
\enditems

\secc Useful formulas

\begitems \style N
\i Factorization of $a^n-b^n$ and $a^n+b^n$ (for odd $n$)
\i Expressions of the form $a^4+4b^4$ can often be factored, even if it is not evident
\i Non-evident factorization of $a^3+b^3+c^3-3abc$ and its consequences (Proposition 3)
\i Inequality $a^2+b^2+c^2 \ge ab+bc+ca$ and its proof
\i Other useful factorizations like 
    \begitems \style i
    \i $a^2+2ab+b^2=(a+b)^2$
    \i $a^2+b^2+c^2+2ab+2bc+2ca=(a+b+c)^2$
    \i etc, there are no limits to the imagination of problem authors...
    \enditems 
\enditems

\DisplayExerciseSolutions

\DisplayHints

\DisplayProblemSolutions

\bye